% Fixture: the machine translation of latex-features-showcase.tex as the build would write it (only the metadata differs, for the tests).
\title{LaTeX-functies in het kort}
\lang{nl}
\date{2026-03-02}
\excerpt{Een compleet overzicht of the LaTeX this blog understands — sectioning, lists, tables, listings, quotes, figures and more.}
\tags{LaTeX, Reference, Web}
\coverimage{/blog-example.png}
\xpost{https://x.com/YOUR_HANDLE/status/PLACEHOLDER_TWEET_ID_1}

\begin{document}

\section{Introduction}

This post demonstrates the LaTeX this blog understands. Use it as a reference for writing your own posts, or just to see how everything renders.

\section{Text Formatting}

You can write \textbf{bold text}, \emph{italic text}, and \emph{\textbf{bold italic}} together. There's also \textasciitilde{}\textasciitilde{}strikethrough\textasciitilde{}\textasciitilde{} for deleted content. Inline \texttt{code} looks like this.

Here is a \href{https://example.com}{regular link} and an auto-linked URL: https://example.com.

\section{Headings}

This post uses \texttt{\#\#} and \texttt{\#\#\#} headings (h2 and h3). They automatically appear in the table of contents on the right.

\subsection{This is an H3 Subheading}

Subheadings help break up longer sections into logical parts.

\section{Lists}

\subsection{Unordered Lists}

\begin{itemize}
  \item First item
  \item Second item with \textbf{bold}
  \item Nested lists work too:
  \item Sub-item one
  \item Sub-item two
  \item Even deeper nesting
  \item Back to top level
\end{itemize}

\subsection{Ordered Lists}

\begin{enumerate}
  \item First step
  \item Second step
  \item Third step
  \item Sub-step A
  \item Sub-step B
  \item Final step
\end{enumerate}

\subsection{Task Lists}

\begin{itemize}
  \item Implement LaTeX rendering
  \item [x] Add syntax highlighting
  \item [ ] Write more blog posts
  \item [ ] World domination
\end{itemize}

\section{Blockquotes}

\begin{quote}[--- Martin Fowler]
``Any fool can write code that a computer can understand. Good programmers write code that humans can understand.''
\end{quote}

The attribution goes in square brackets after \texttt{\textbackslash{}begin\{quote\}} and is set flush right underneath.

Nested blockquotes also work:

\begin{quote}
This is the outer quote.

\begin{quote}
And this is a nested quote inside it.
\end{quote}

Back to the outer quote.
\end{quote}

\section{Code Blocks}

Inline code: \texttt{const x = 42;}

Fenced code block with syntax highlighting:

\begin{lstlisting}[language=typescript]
interface BlogPost {
  slug: string;
  title: string;
  date: string;
  tags: string[];
}

function getLatestPosts(posts: BlogPost[], count: number): BlogPost[] {
  return posts
    .sort((a, b) => (a.date > b.date ? -1 : 1))
    .slice(0, count);
}
\end{lstlisting}

A Python example:

\begin{lstlisting}[language=python]
def fibonacci(n: int) -> list[int]:
    """Generate the first n Fibonacci numbers."""
    fib = [0, 1]
    for i in range(2, n):
        fib.append(fib[i - 1] + fib[i - 2])
    return fib[:n]

print(fibonacci(10))
# [0, 1, 1, 2, 3, 5, 8, 13, 21, 34]
\end{lstlisting}

A shell command:

\begin{lstlisting}[language=bash]
git log --oneline --graph --all | head -20
\end{lstlisting}

A listing can carry a file name in its header — \texttt{title=} for \texttt{lstlisting}, \texttt{filename=} for \texttt{minted}:

\begin{minted}[filename=blink.py]{python}
from machine import Pin
import time

led = Pin(25, Pin.OUT)
while True:
    led.toggle()
    time.sleep(0.5)
\end{minted}

And \texttt{verbatim} prints text exactly as written, without colouring:

\begin{verbatim}
$ make flash
avrdude: 1024 bytes of flash written
\end{verbatim}

\section{Tables}

\begin{tabular}{lcl}
\hline
Feature & Supported & Notes \\
\hline
Bold & ✅ & \texttt{**text**} \\
Italic & ✅ & \texttt{*text*} \\
Strikethrough & ✅ & \texttt{~~text~~} \\
Tables & ✅ & GFM extension \\
Task lists & ✅ & GFM extension \\
Footnotes & ✅ & GFM extension \\
\hline
\end{tabular}

\subsection{Alignment in Tables}

\begin{tabular}{lcr}
\hline
Left-aligned & Center-aligned & Right-aligned \\
\hline
Content & Content & Content \\
More & More & More \\
Data & Data & Data \\
\hline
\end{tabular}

\section{Footnotes}

Footnotes use \texttt{\textbackslash{}footnote\{\}}.\footnote{The text of the footnote, collected at the bottom of the page.} You can also write longer ones.\footnote{A longer footnote can hold \textbf{bold}, \texttt{code} and math such as $e^{i\pi} + 1 = 0$.}

\section{Images}

Figures use the \texttt{figure} environment. The \texttt{\textbackslash{}caption\{\}} becomes the caption, and each figure gets an automatic number you can point at once you give it a \texttt{\textbackslash{}label\{\}}: \texttt{\textbackslash{}Cref\{\}} prints the word and the number — this one is \Cref{fig-example} — and \texttt{\textbackslash{}ref\{\}} the number alone, exactly as LaTeX does.

\begin{figure}
\includegraphics[width=0.4\textwidth]{/blog-example.png}
\caption{Example image}
\label{fig-example}
\end{figure}

As shown in \href{#fig-1}{Fig. 1}, images are automatically numbered and captioned.

\section{Diagrams}

A listing in the language \texttt{mermaid} is drawn as a diagram when the site is built; \Cref{fig-flow} is Mermaid text. Put it in a \texttt{figure} for a number and a caption, and give it a \texttt{title=} for screen readers.

\begin{figure}
\begin{lstlisting}[language=mermaid, title={The control loop}]
flowchart LR
  A[Sensor] --> B[Controller]
  B --> C{In range?}
  C -->|yes| D[Actuator]
  C -->|no| E[Alarm]
\end{lstlisting}
\caption{A flowchart drawn from Mermaid text}
\label{fig-flow}
\end{figure}

\section{Emphasis Combinations}

Here are some formatting combos:

\begin{itemize}
  \item \textbf{Bold with \texttt{inline code} inside}
  \item \emph{Italic with \href{https://example.com}{a link} inside}
  \item \textasciitilde{}\textasciitilde{}Strikethrough with \textbf{bold} inside\textasciitilde{}\textasciitilde{}
  \item A paragraph with \textbf{bold}, \emph{italic}, \texttt{code}, and a \href{https://example.com}{link} all mixed together
\end{itemize}

\section{Long Code Block}

To test scrolling behavior with longer blocks:

\begin{lstlisting}[language=c]
#include <stdio.h>
#include <stdlib.h>

typedef struct Node {
    int data;
    struct Node* next;
} Node;

Node* create_node(int data) {
    Node* node = malloc(sizeof(Node));
    node->data = data;
    node->next = NULL;
    return node;
}

void insert_front(Node** head, int data) {
    Node* new_node = create_node(data);
    new_node->next = *head;
    *head = new_node;
}

void print_list(Node* head) {
    Node* current = head;
    while (current != NULL) {
        printf("%d -> ", current->data);
        current = current->next;
    }
    printf("NULL\n");
}

int main(void) {
    Node* list = NULL;
    for (int i = 5; i >= 1; i--) {
        insert_front(&list, i);
    }
    print_list(list); // 1 -> 2 -> 3 -> 4 -> 5 -> NULL
    return 0;
}
\end{lstlisting}

\section{Line Breaks \& Paragraphs}

This is the first paragraph. LaTeX needs a blank line between paragraphs to separate them.

This is the second paragraph. Without a blank line, text would continue on the same paragraph.

\section{Summary}

That covers the structural LaTeX: text formatting, sectioning, lists, quotes, listings, tables and figures. Everything here is parsed with \texttt{@unified-latex} and rendered to HTML at build time.

\end{document}

